% work_area: work as the area under the F-x graph (Ch 11, fig:workarea)
% A spring gun: the spring starts compressed by x0 and pushes the ball while it
% relaxes.  After the ball has moved x the compression is x0 - x, so
%   F = k(x0 - x),  falling linearly from kx0 to 0 at x = x0.
% One thin strip (height F at its midpoint) has area F dx = dW; the strips add to the triangle
%   W = (1/2)(kx0)(x0) = (1/2) k x0^2,  the energy the spring stored (Sec. on U).
% Units: k = 1, x0 = 1.
\begin{tikzpicture}
\begin{axis}[bookplot,
    width=9.4cm, height=5.6cm,
    xmin=0, xmax=1.22, ymin=0, ymax=1.2,
    xlabel={distance moved $x$}, ylabel={force $F$},
    ylabel style={at={(axis description cs:0,1)},anchor=south,rotate=-90,yshift=1pt},
    xtick={1}, xticklabels={$x_0$},
    ytick={1}, yticklabels={$kx_0$}]
  % the area: the work done on the ball
  \addplot[area fill,domain=0:1,samples=2] {1-x} \closedcycle;
  % one thin strip, F dx
  \def\xs{0.56}\def\dx{0.06}
  \draw[strip,fill=cPrim!40] (axis cs:\xs,0) rectangle (axis cs:\xs+\dx,{1-\xs-0.5*\dx});
  % the force line
  \addplot[cPrim,domain=0:1,samples=2] {1-x};
  \node[cPrim,anchor=south west] at (axis cs:0.30,0.72) {$F=k(x_0-x)$};
  % labels
  \draw[leader,cPrim] (axis cs:\xs+0.5*\dx,0.30) -- (axis cs:0.80,0.55)
    node[anchor=south west,inner sep=1pt] {strip: $dW=F\,dx$};
  \node[cPrim,align=left] at (axis cs:0.25,0.25) {$W=\text{area}$\\[1pt] $\phantom{W}=\tfrac12kx_0^2$};
\end{axis}
\end{tikzpicture}
